% Part: computability
% Chapter: recursive-functions
% Section: pr-functions

\documentclass[../../../include/open-logic-section]{subfiles}

\begin{document}

\olfileid{cmp}{rec}{prf}
\olsection{بنیادی بازگشتی فنکشن}

آؤ فیر درج کر لَیئے کہ موجودہ فنکشناں توں
بنیادی بازگشت تے ترکیب ورت کے نویں فنکشن
کِویں تعریف کیتے جا سکدے نیں۔

\begin{defn}
  \ollabel{defn:primitive-recursion} فرض کرو $f$ اک $k$-جگہ
  فنکشن اے ($k\ge 1$) تے $g$ اک $(k+2)$-جگہ فنکشن اے۔
  $f$ تے $g$ توں \emph{بنیادی بازگشت راہیں متعین کیتا فنکشن}
  اوہ $(k+1)$-جگہ فنکشن~$h$ اے جیہڑا ایہناں مساواتاں نال
  متعین ہوندا اے:
  \begin{align*}
    h(x_0,\dots,x_{k-1},0) & =  f(x_0,\dots,x_{k-1}) \\
    h(x_0,\dots,x_{k-1},y+1) & = g(x_0,\dots,x_{k-1}, y, h(x_0,\dots,x_{k-1}, y))
  \end{align*}
\end{defn}

\begin{defn}
  \ollabel{defn:composition} فرض کرو $f$ اک $k$-جگہ فنکشن اے،
  تے $g_0$، \dots، $g_{k-1}$ کل $k$ فنکشن نیں، جیہڑے
  سارے $n$-جگہ نیں۔ $f$ تے
  \emph{$g_0$، \dots،~$g_{k-1}$ توں ترکیب راہیں متعین کیتا فنکشن}
  اوہ $n$-جگہ فنکشن~$h$ اے جیہڑا ایویں متعین ہوندا اے:
  \[
  h(x_0, \dots, x_{n-1}) =
  f(g_0(x_0, \dots, x_{n-1}), \dots, g_{k-1}(x_0, \dots, x_{n-1})).
  \]
\end{defn}

$\Succ$ تے پروجیکشن فنکشناں
\[
\Proj{n}{i}(x_0,\dots,x_{n-1}) = x_i,
\]
دے نال، ہر قدرتی عدد $n$ تے $i < n$ لئی،
اسی فنکشن~$\Zero(x) = 0$ نوں وی بنیادی
بازگشتی فنکشناں وچ شامل کراں گے۔

\begin{defn}
  بنیادی بازگشتی فنکشناں دا سیٹ، $\Nat^n$ توں
  $\Nat$ ول فنکشناں دا اوہ سیٹ اے جیہڑا
  ہیٹھ لکھیاں شرطاں نال استقرائی طور اُتے
  متعین ہوندا اے:
  \begin{enumerate}
  \item $\Zero$ بنیادی بازگشتی اے۔
  \item $\Succ$ بنیادی بازگشتی اے۔
  \item ہر پروجیکشن فنکشن $\Proj{n}{i}$ بنیادی بازگشتی اے۔
  \item جے $f$ اک $k$-جگہ بنیادی بازگشتی فنکشن اے
    تے $g_0$، \dots،~$g_{k-1}$، $n$-جگہ
    بنیادی بازگشتی فنکشن نیں، تاں $f$ دی
    $g_0$، \dots،~$g_{k-1}$ نال ترکیب وی
    بنیادی بازگشتی اے۔
  \item جے $f$ اک $k$-جگہ بنیادی بازگشتی فنکشن اے
    تے $g$ اک $k+2$-جگہ بنیادی بازگشتی فنکشن اے،
    تاں $f$ تے $g$ توں بنیادی بازگشت راہیں
    متعین کیتا فنکشن وی بنیادی بازگشتی اے۔
\end{enumerate}
\end{defn}

\begin{explain}
مختصر لفظاں وچ، بنیادی بازگشتی فنکشناں دا سیٹ
اوہ سب توں چھوٹا سیٹ اے جیہڑا $\Zero$،
$\Succ$ تے پروجیکشن فنکشناں~$\Proj{n}{j}$
نوں رکھدا اے، تے ترکیب تے بنیادی بازگشت
تحت بند اے۔

بنیادی بازگشتی فنکشناں دے سیٹ نوں بیان کرن
دا اک ہور طریقہ اوہنوں ``مرحلیاں'' دے
حوالے نال متعین کرنا اے۔ $S_0$ شروعاتی
فنکشناں دا سیٹ ہووے: $\Zero$، $\Succ$
تے پروجیکشن فنکشن۔ ایہ مرحلے~$0$ دے
بنیادی بازگشتی فنکشن نیں۔ جے مرحلہ $S_i$
متعین ہو چکیا ہووے، تاں $S_{i+1}$ اوہناں
سارے فنکشناں دا سیٹ ہووے جیہڑے پہلے توں~$S_i$
وچ موجود فنکشناں اُتے ترکیب یا بنیادی بازگشت
دی اکّو مثال لا کے ملدے نیں۔ فیر
\[
S = \bigcup_{i \in \Nat} S_i
\]
سارے بنیادی بازگشتی فنکشناں دا سیٹ اے۔
\end{explain}

آؤ تصدیق کریے کہ $\Add$ بنیادی بازگشتی
فنکشن اے۔

\begin{prop}
  جمع دا فنکشن $\Add(x,y) = x+y$ بنیادی بازگشتی اے۔
\end{prop}

\begin{proof}
ساڈے کول $\Add$ دی بنیادی بازگشتی تعریف
پہلے ہی دو فنکشناں $f$ تے~$g$ دے
ذریعے اے، جیہڑی \olref{defn:primitive-recursion}
دے نمونے نال ملدی اے:
\begin{align*}
  \Add(x_0, 0) & = f(x_0) = x_0\\
  \Add(x_0, y+1) & = g(x_0, y, \Add(x_0, y)) =
  \Succ(\Add(x_0, y))
\end{align*}
ایس لئی $\Add$ بنیادی بازگشتی اے،
بشرطیکہ $f$ تے $g$ وی بنیادی بازگشتی ہون۔
$f(x_0) = x_0 = \Proj{1}{0}(x_0)$،
تے پروجیکشن فنکشن بنیادی بازگشتی گنے جاندے
نیں، ایس لئی $f$ بنیادی بازگشتی اے۔
فنکشن $g$ اوہ تِن-جگہ فنکشن
$g(x_0, y, z)$ اے جیہڑا ایویں متعین اے:
\[
g(x_0, y, z) = \Succ(z).
\]
ایس توں ہن تک ایہ نتیجہ نہیں نکلدا کہ
$g$ بنیادی بازگشتی اے، کیوں جے $g$
تے $\Succ$ بالکل اکّو فنکشن نہیں:
$\Succ$ اک-جگہ اے تے $g$ نوں تِن-جگہ
ہونا پیندا اے۔ پر اسی $g$ نوں ترکیب
راہیں باقاعدہ ایویں تعریف کر سکدے آں:
\[
g(x_0, y, z) = \Succ(\Proj{3}{2}(x_0, y, z))
\]
چونکہ $\Succ$ تے $\Proj{3}{2}$
بنیادی بازگشتی فنکشن گنے جاندے نیں،
ایس لئی $g$ وی بنیادی بازگشتی اے،
کیونکہ اوہ بنیادی بازگشتی فنکشناں توں
ترکیب راہیں متعین ہو سکدا اے۔
\end{proof}

\begin{prop}
  \ollabel{prop:mult-pr}
  ضرب دا فنکشن $\Mult(x,y) = x \cdot y$ بنیادی بازگشتی اے۔
\end{prop}

\begin{proof}
  مشق۔
\end{proof}

\begin{prob}
  \olref[cmp][rec][prf]{prop:mult-pr} ثابت کرو:
  دکھاؤ کہ $\Mult$ دی بنیادی بازگشتی تعریف
  نوں \olref[cmp][rec][prf]{defn:primitive-recursion}
  دے مطلوبہ روپ وچ لکھیا جا سکدا اے، تے
  متعلقہ فنکشن $f$ تے~$g$ بنیادی بازگشتی نیں۔
\end{prob}

\begin{ex}
ایہ ساڈی بنیادی بازگشتی تعریف دی سب توں پہلی
مثال اے:
\begin{align*}
h(0) & =  1 \\
h(y+1) & =  2 \cdot h(y).
\end{align*}
ایہ فنکشن \olref{defn:primitive-recursion} دے
مطلوبہ روپ وچ نہیں آ سکدا، کیونکہ $k=0$۔
تعریف وچ مستقل عدد $1$ تے $2$ وی آؤندے
نیں۔ پہلے مسئلے نوں حل کرن لئی اک اضافی
غیر مستعمل دلیل شامل کریے تے فنکشن~$h'$
ایویں تعریف کریے:
\begin{align*}
h'(x_0, 0) & =  f(x_0) = 1 \\
h'(x_0, y+1) & =  g(x_0, y, h'(x_0, y)) = 2 \cdot h'(x_0, y).
\end{align*}
فنکشن $f(x_0) = 1$ نوں $\Succ$
تے $\Zero$ دی ترکیب نال تعریف کر
سکدے آں: $f(x_0) = \Succ(\Zero(x_0))$۔
فنکشن $g$ نوں $g'(z) = 2 \cdot z$
تے پروجیکشن فنکشناں توں ترکیب راہیں
تعریف کر سکدے آں:
\begin{align*}
g(x_0, y, z) & = g'(\Proj{3}{2}(x_0, y, z))
\intertext{تے $g'$ نوں اگے ترکیب راہیں ایویں تعریف کر سکدے آں}
g'(z) & = \Mult(g''(z), \Proj{1}{0}(z))
\intertext{تے}
g''(z) & = \Succ(f(z)),
\end{align*}
جتھے $f$ اوہی اُپر والا اے:
$f(z) = \Succ(\Zero(z))$۔ ہن جدوں
ساڈے کول~$h'$ اے، تاں اسی فیر ترکیب
ورّت کے $h(y) =
h'(\Proj{1}{0}(y),\Proj{1}{0}(y))$
رکھ سکدے آں۔ ایہ دکھاؤندا اے کہ $h$
نوں بنیادی فنکشناں توں ترکیباں تے
بنیادی بازگشتاں دی اک ترتیب نال تعریف
کیتا جا سکدا اے، ایس لئی $h$~بنیادی
بازگشتی اے۔
\end{ex}

\end{document}
